\documentclass[12pt,a4paper,answers]{exam}
\usepackage[T1]{fontenc}
\usepackage{longtable}
\usepackage{graphicx}

\usepackage{hyperref}
%\urlstyle{same}  % don't use monospace font for urls
\setlength{\parindent}{0pt}
\setlength{\parskip}{6pt plus 2pt minus 1pt}

\usepackage{geometry}
\usepackage{multicol}
\usepackage{relsize}
\usepackage{pgf}
\usepackage{minted}

\geometry{
  %includeheadfoot,
  margin=2.54cm
}
\title{Exercises on Design Patterns V}
\author{Last one\ldots}
\date{2018}

\fvset{tabsize=2}
%\fvset{fontsize=\relsize{-1}}
%\fvset{frame=single}

% https://www.javacodegeeks.com/2015/09/decorator-design-pattern.html

\newcommand{\code}{/Users/keith/Documents/Courses/sdp/repo/worksheets/dp-five/scala/src}
%\newcommand{\code}{/Users/keith/Courses/sdp/SDP-2017/exercises/week11/src/main/scala}
\newcommand{\codesol}{/Users/keith/Courses/sdp/SDP-2017/exercises/week11sol/src/main/scala}

\begin{document}
\maketitle

In these exercises we will be examining the following design patterns:
\begin{itemize}
\item State,
\item Template Method, and
\item Visitor.	
\end{itemize}


You should use Scala or Kotlin to answer the exercises.\\
The source code examples can be found on the repository under the 
\verb!worksheets/dp-five! folder.

\begin{questions}

%\inputminted{scala}{\codesol/composite/HtmlElement.scala}

\question
This question concerns the \textsc{State} design pattern.

To illustrate the use of the \textsc{State Design Pattern}, you will help a company which is 
looking to build a robot for cooking. 
The company wants a simple robot that can simply \emph{walk} and \emph{cook}. 
A user can operate a robot using a set of commands via remote control. 
Currently, a robot can do three things, it can \emph{walk}, \emph{cook}, or can be switched off.

The company has set protocols to define the functionality of the robot. 
If a robot is in an \textbf{on} state you can command it to walk. 
If asked to cook, the state would change to \textbf{cook} or if set to \textbf{off}, it will be 
switched off.

Similarly, when in \textbf{cook} state it can walk or cook, but cannot be switched off. And 
finally, when in \textbf{off} state it will automatically turn on and walk when the user commands 
it to walk but cannot cook in the off state.

This might look like an easy implementation: a robot class with a set of methods like 
\mintinline{scala}!walk!, 
\mintinline{scala}!cook!, \mintinline{scala}!off!, and states like \textbf{on}, \textbf{cook}, and 
\textbf{off}. 
We can use \emph{if-else} branches, or switches, to implement the protocols set by the company. 
Too much \emph{if-else} or switch statements will create a maintenance nightmare as complexity 
might increase in the future.

You might think that we can implement this without issues using \emph{if-else} statements, but as 
a change comes the code would become more complex. 
The requirement clearly shows that the behaviour of an object is truly based on the state of that 
object. 
To achieve this you can use the \textsc{State Design Pattern} which encapsulates the states of 
the object into another individual class and keeps the context class independent of any state 
change.

The following is the \mintinline{scala}!RoboticState! interface which contains the behaviour of a 
robot.
\inputminted[]{scala}{\code/state/RoboticState.scala}
The \mintinline{scala}!Robot! class is a concrete class implements the 
\mintinline{scala}!RoboticState! interface. 
The class contains the set of all possible states a robot can be in.
\inputminted[]{scala}{\code/state/Robot.scala}
The class initialises all the states and sets the current state as on.
Now, we will see all the concrete states of a robot. 
A robot will be in any of these states at any time.
\inputminted{scala}{\code/state/RoboticOn.scala}
\inputminted{scala}{\code/state/RoboticCook.scala}
\inputminted{scala}{\code/state/RoboticOff.scala}
Now, let's test the code.
\inputminted{scala}{\code/state/TestStatePattern.scala}
The above code should result in the following output:
\begin{Verbatim}
Walking...
Cooking...
Walking...
Robot is switched off
Walking...
Robot is switched off
Cannot cook at Off state.	
\end{Verbatim}

\begin{solution}
	\ldots
\end{solution}

\question
The \textsc{Template Design Pattern} is a behaviour pattern and, as the name suggests, 
it provides a template or a structure of an algorithm which is used by users. 
A user provides its own implementation without changing the algorithm's structure.

Have you ever connected to a relation database using your Java application? 
Let's recall some important steps which are required to connect and insert data into 
the database. First, we need a driver according to the database we want to connect with. 
Then, we pass some credentials to the database, then, we prepare a statement, set data into the 
insert statement and insert it using the insert command. Later, we close all the connections, and 
optionally destroy all the connection objects.

You need to write all these steps regardless of any vendor's relational database. 
Consider a problem where you need to insert some data into the different databases. You need to 
fetch some data from a CSV file and have to insert it into a MySQL database. Some data comes from 
a text file and which should be insert into an Oracle database. The only difference is the driver 
and the data, the rest of the steps should be the same, as JDBC provides a common set of 
interfaces to communicate to any vendor's specific relation database.

We can create a template, which will perform some steps for the client, and we will leave 
some steps to let the client to implement them in its own specific way. 
Optionally, a client can override the default behaviour of some already defined steps.

Below we can see the connection template class used to provide a template 
for clients to connect and communicate with the various databases.
\inputminted{scala}{\code/template/ConnectionTemplate.scala}
The abstract class provides steps to connect, communicate and later to close the connections. All 
these steps are required in order to get the work done. The class provides default implementation 
to some common steps and leaves the specific steps as abstract which force the client to provide 
an implementation to them.

The \mintinline{scala}!setDBDriver! method should be implementing by the user to provide 
database specific drivers. The credentials could be different for different databases; 
therefore, \mintinline{scala}!setCredentials! is also left as abstract to let the 
user implement it.

Similarly, connecting to the database using the JDBC API and preparing a statement is common.
But, data would be specific so user will provide it, and rest of other steps like running an
insert statement, closing connections and destroying object are similar to any 
database, therefore their implementations are kept common inside the template.

The key method of the above class is the \mintinline{scala}!run! method. 
The \mintinline{scala}!run! method is used to run these steps in order. 
The method is set as final because the steps should be kept safe and should not be change 
by any user.

The below two classes extend the template class and provide specific 
implementation to some methods.
\inputminted{scala}{\code/template/MySqLCSVCon.scala}
The above class is used to connect to a MySQL database and provides data by reading a CSV file.
\inputminted{scala}{\code/template/OracleTxtCon.scala}
The above class is used to connect to an Oracle database and provides 
data by reading a text file.
\inputminted{scala}{\code/template/TestTemplatePattern.scala}
The above code should result in the following output:
\begin{Verbatim}
For MYSQL....
Setting MySQL DB drivers...
Setting credentials for MySQL DB...
Setting connection...
Preparing insert statement...
Setting up data from csv file....
Inserting data...
Closing connections...
Destroying connection objects...

For Oracle...
Setting Oracle DB drivers...
Setting credentials for Oracle DB...
Setting connection...
Preparing insert statement...
Setting up data from txt file....
Inserting data...
Closing connections...
Destroying connection objects...	
\end{Verbatim}
The above output clearly shows how the template pattern works to connect and communicate 
with the different databases using the similar way. The pattern keeps the common code under 
one class and promotes code reusability. It sets a framework and controls it for the users 
and allows the users to extends the template in order to provide their specific 
implementation to some of the steps.
\begin{solution}
	\ldots
\end{solution}

\question
Implementing the \textsc{Visitor Design Pattern} requires two important interfaces, 
an \mintinline{scala}!Element! interface which will contain an 
\mintinline{scala}!accept! method with an argument of type 
\mintinline{scala}!Visitor!. This interface will be implemented by all the classes that need 
to allow visitors to visit them. In our case, the class \mintinline{scala}!HtmlTag! 
will implement the \mintinline{scala}!Element! interface, as the 
\mintinline{scala}!HtmlTag! is the parent abstract class of all the concrete html classes, 
the concrete classes will inherit and will override the \mintinline{scala}!accept! 
method of the \mintinline{scala}!Element! interface.

The other important interface is the \mintinline{scala}!Visitor! interface; this 
interface will contain \mintinline{scala}!visit! methods with an argument of a 
class that implements the \mintinline{scala}!Element! interface. 
Please also note that we have added two new methods in our 
\mintinline{scala}!HtmlTag! class, the \mintinline{scala}!getStartTag! and
\mintinline{scala}!getEndTag!. 
\inputminted{scala}{\code/visitor/Element.scala}
and
\inputminted{scala}{\code/visitor/Visitor.scala}
The supplemental classes are:
\inputminted{scala}{\code/visitor/HtmlTag.scala}
The abstract \mintinline{scala}!HtmlTag! class implements the \mintinline{scala}!Element!
interface. The below concrete classes will override the 
\mintinline{scala}!accept! method of the \mintinline{scala}!Element! interface 
and will call the \mintinline{scala}!visit! method, and will pass this operator as an argument.
This will allow the visitor method to get all the public members of the object, 
and to add new operations based on it.
\inputminted{scala}{\code/visitor/HtmlParentElement.scala}
\inputminted{scala}{\code/visitor/HtmlElement.scala}
Now, the concrete visitor classes: we have created two concrete classes, 
one will add a css class visitor to all html tags and the other one will 
change the width of the tag using the style attribute of the html tag.
\inputminted{scala}{\code/visitor/CssClassVisitor.scala}
Now, let's test your code.
\inputminted{scala}{\code/visitor/TestVisitorPattern.scala}
The above code should result in the following output:
\begin{Verbatim}
Before visitor.........

<html>
<body>
<P>Testing html tag library</P>
<P>Paragraph 2</P>
</body>
</html>

After visitor.......

<html style='width:58px;' class='visitor'>
<body style='width:58px;' class='visitor'>
<p style='width:46px;' class='visitor'>Testing html tag library</P>
<p style='width:46px;' class='visitor'>Paragraph 2</P>
</body>
</html>	
\end{Verbatim}

\begin{solution}
	\ldots
\end{solution}
\end{questions}

\end{document}